% !TeX root = ../../Book5.tex
% 纲要标题：### 6.3 Garnir 关系与标准基
% 纲要位置：工作记录/计划纲要.md，第 3919 行。
% * 建立 Garnir 关系和标准化过程，证明标准列交错元的生成性及线性无关性
% * 固定表的次序和标准化规则，分别说明重写为何终止、标准列交错元为何生成，以及一般系数环上哪些步骤仍成立。
% * 以首项和三角性组织线性无关性证明，区分整系数标准基与基变换后的域上模型；不要仅凭标准表个数推断线性无关。
\section{Garnir 关系与标准基}
\label{b5:sec:0603}

本节先在整数系数下处理关系。这样既能避免用除法掩盖特征限制，也能同时得到任意域和任意交换环上的标准基。整个过程只移动表中数字，不改变 Young 图的形状。

\subsection{Garnir 关系与标准列交错元基}
\label{b5:subsec:0603:01}

\begin{lemma}[Garnir 关系]\label{b5:lem:garnir}
设 \(n\) 为非负整数，\(\lambda\vdash n\)，\(\lambda'\) 为共轭分拆，\(t\) 为形状 \(\lambda\) 的 Young 表。在相邻的第 \(j,j+1\) 列中，分别选数字子集 \(A,B\)，并假定
\[
|A|+|B|>\lambda'_j.
\]
令 \(H=S_A\times S_B\leq S_{A\cup B}\)，选取左陪集 \(S_{A\cup B}/H\) 的代表元集 \(D\)，其中包含单位元。记 \(e_u\) 为表 \(u\) 的列交错元，则在整数行表类置换模 \(M_{\mathbb Z}^\lambda\) 中有
\begin{equation}\label{b5:eq:garnir}
\sum_{\sigma\in D}\operatorname{sgn}(\sigma)e_{\sigma t}=0.
\end{equation}
同一等式经系数映射后，在任意交换含幺环上仍成立。它称为 \term[Garnir-guanxi]{Garnir 关系}[Garnir relation]。
\end{lemma}
\begin{proof}
记 \(\gamma=\sum_{\pi\in S_{A\cup B}}\operatorname{sgn}(\pi)\pi\)。对任意 \(c\in C_t\)，数字集合 \(A\cup B\) 在 \(ct\) 中仍全部位于这两列，因而位于前 \(\lambda'_j\) 行。由所假定的不等式，其中必有两个数字处于同一行；交换这两个数字的换位 \(\tau\in S_{A\cup B}\) 固定 \(\{ct\}\)。在 \(\gamma\{ct\}\) 的求和中把 \(\pi\) 与 \(\pi\tau\) 配成一对，两项符号相反、向量相同，故它们相消。因此 \(\gamma e_t=0\)。

另一方面，按左陪集分组得
\[
\gamma=\sum_{\sigma\in D}\operatorname{sgn}(\sigma)\sigma
\sum_{h\in H}\operatorname{sgn}(h)h.
\]
\(H\) 包含在 \(C_t\) 中，故内层和作用于 \(e_t\) 时得到 \(|A|!\,|B|!\,e_t\)。于是
\[
|A|!\,|B|!\sum_{\sigma\in D}\operatorname{sgn}(\sigma)e_{\sigma t}=0.
\]
自由交换群 \(M_{\mathbb Z}^\lambda\) 无挠，故可消去这个非零整数，得到\eqref{b5:eq:garnir}。这是整数系数向量恒等式，随后可映到任意 \(R\)；这一顺序并未在 \(R\) 中除以阶乘。
\end{proof}

\begin{theorem}[标准基]\label{b5:thm:specht-standard-basis}
设 \(R\) 为交换含幺环，\(n\) 为非负整数，\(\lambda\vdash n\)。所有形状为 \(\lambda\) 的标准表 \(t\) 所给的列交错元 \(e_t\)，构成 Specht 模 \(S_R^\lambda\) 的一组 \(R\)-基。因而 \(S_{\mathbb Z}^\lambda\) 是有限秩自由交换群，且自然映射
\[
R\otimes_{\mathbb Z}S_{\mathbb Z}^\lambda\longrightarrow S_R^\lambda
\]
为同构。
\end{theorem}
\begin{proof}
先证明生成性。利用\cref{b5:prop:specht-column-sign}，可先把每列排成递增次序。对列标准表 \(t\)，把第一列从上到下的数字、第二列从上到下的数字依次接起来，所得长度为 \(n\) 的列字记为 \(w(t)\)，按通常字典序比较。可能的列字只有有限多个。

若 \(t\) 还不是标准表，则某一行有相邻逆序
\(t(i,j)>t(i,j+1)\)。取
\[
\begin{aligned}
A&=\Bigl\{t(i,j),t(i+1,j),\ldots,t(\lambda'_j,j)\Bigr\},\\
B&=\Bigl\{t(1,j+1),\ldots,t(i,j+1)\Bigr\}.
\end{aligned}
\]
由于两列递增，\(A\) 中每个数字都大于 \(B\) 中每个数字，且
\(|A|+|B|=\lambda'_j+1\)。在 Garnir 关系中单独移出单位元项，得到
\begin{equation}\label{b5:eq:garnir-rewrite}
e_t=-\sum_{\substack{\sigma\in D\\\sigma\ne1}}
\operatorname{sgn}(\sigma)e_{\sigma t}.
\end{equation}
一个非单位陪集满足 \(\sigma A\ne A\)，所以新的第 \(j\) 列从 \(A\) 中移出若干数字，并换入同样多个来自 \(B\) 的较小数字；前 \(j-1\) 列不变。将新表的每列重新排序后，列字严格小于 \(w(t)\)。列排序只改变列交错元的符号。因此重复\eqref{b5:eq:garnir-rewrite}时，每一条重写链的列字都严格下降；有限性保证过程终止。终止项不能还有相邻行逆序，所以全部为标准表，生成性得证。

再证明线性无关。对任意行表类 \(\{u\}\)，令
\[
N_u(a,m)=\#\{\text{前 }a\text{ 行中不大于 }m\text{ 的数字}\}.
\]
这只依赖表类。若 \(t\) 标准，则对每个 \(c\in C_t\) 都有
\begin{equation}\label{b5:eq:tabloid-triangular}
N_{ct}(a,m)\leq N_t(a,m)\qquad\text{对所有 }a,m.
\end{equation}
事实上，第 \(j\) 列中不大于 \(m\) 的数字在 \(t\) 中占据最上面的若干位置，若其个数为 \(q_j\)，则前 \(a\) 行贡献 \(\min(a,q_j)\)；列置换以后的贡献不会更大。对各列相加即得不等式。

在标准表集合上，用所有这些不等式定义偏序。若双向成立，则每一行中小于等于每个 \(m\) 的数字个数都相同，因此各行的数字集合相同；标准表各行递增，故两个表相同。把这个有限偏序延伸为全序。在 \(e_t\) 的表类展开中，若标准表 \(u\) 的表类 \(\{u\}\) 出现，则\eqref{b5:eq:tabloid-triangular}说明 \(u\) 不大于 \(t\)；而 \(\{t\}\) 的系数由\cref{b5:prop:specht-column-sign}知为 \(1\)。于是只取标准表所对应的那些表类坐标，得到一个对角元全为 \(1\) 的三角矩阵。它在任意交换环上可逆，所以这些 \(e_t\) 线性无关。

最后，生成性所用系数全为整数；自然映射把整数标准基变成 \(R\) 上刚证的标准基，故为同构。
\end{proof}

\subsection{表的次序、标准化终止与生成性}
\label{b5:subsec:0603:02}

上述证明给出实际的标准化算法。输入为一个填数表，或若干列交错元的有限 \(R\)-线性组合；输出为标准列交错元的线性组合。每次先排序各列并记下符号，再找到一处相邻行逆序，取证明中规定的两段数字，按\eqref{b5:eq:garnir-rewrite}展开，最后合并相同标准表的系数。这个算法只用加减法，不要求能在 \(R\) 中求逆。

\begin{example}\label{b5:exm:garnir-21}
对列标准而非标准的表
\[
t=\begin{array}{cc}2&1\\3&\end{array}
\]
取 \(A=\{2,3\}\)、\(B=\{1\}\)，左陪集代表可取 \(1,(12),(13)\)。Garnir 关系为
\(e_t-e_{(12)t}-e_{(13)t}=0\)。令
\[
u=\begin{array}{cc}1&2\\3&\end{array},
\qquad
v=\begin{array}{cc}1&3\\2&\end{array}.
\]
则 \((12)t=u\)，而 \((13)t\) 在第一列交换上下数字后变成 \(v\)，所以
\[
e_t=e_u-e_v.
\]
在\cref{b5:exm:specht-21}的记号下，这正是
\(v_3-v_2=(v_3-v_1)-(v_2-v_1)\)。
\end{example}

算法终止并不自动保证计算很快。每次 Garnir 展开可能产生多个新项，本章只证明有限步骤内能完成重写，不给出随 \(n\) 增长的复杂度界。

\subsection{首项三角性与整系数线性无关性}
\label{b5:subsec:0603:03}

在独立性证明中，起作用的不仅是“存在一个首项”，还包括首项系数是单位 \(1\)。若这个系数只是某个非零整数，约化到正特征后可能消失；这里的三角矩阵在所有系数环上仍可逆。因此标准基不因特征改变而丢失，但由它给出的表示却可能变得可约。

\ProseParagraphBreak

记形状 \(\lambda\) 的标准表数为 \(f^\lambda\)。标准基定理现在才允许我们写
\[
\operatorname{rank}_{\mathbb Z}S_{\mathbb Z}^\lambda=f^\lambda,
\qquad
\dim_kS_k^\lambda=f^\lambda
\]
对任意域 \(k\) 成立。后面计算 \(f^\lambda\) 的公式用于求维数；标准基的存在已经证明，不可约分类将另作论证。
